\documentclass{book}
\usepackage{ctex}
\usepackage[cmex10]{amsmath}
\usepackage{xcolor}
\usepackage{mathtools}
\usepackage{booktabs}
\DeclareRobustCommand{\bluebold}[1]{\textbf{\textcolor[rgb]{0.31,0.63,0.86}{#1}}}
\usepackage[
    pdfpagemode=UseOutlines,
    bookmarks=true,
    bookmarksopen=true,
    colorlinks=true,
    linkcolor=black,
    citecolor=black,
    urlcolor=black
]{hyperref}
\usepackage{bookmark}

\raggedbottom

\begin{document}
\tableofcontents
\clearpage


流体流动现象知识点

\begin{table}[htbp]
\centering
\begin{tabular}{c c} 
\toprule
知识点 & point\\ [0.5ex] 
\midrule
 1 & 流量与流速 \\ 
 2 & 稳态流动与非稳态流动 \\
 3 & 牛顿黏性定律与流体黏性 \\
 4 & 流体类型与雷诺数 \\
 5 & 流体在圆管内的速度分布 \\ [1ex] 
 5 & 流动边界层的概念 \\ [1ex] 
\bottomrule
\end{tabular}
\caption{Knowledge Point Table of Fluid Flow Phenomena.}
\label{table:data}
\end{table}

\chapter{流体流动现象}

\section{流量}

在单位时间内流过管道内任一截面流体的量。
体积流量 $\mathrm{q_v}$，$\mathrm{m^3/s}$ 和质量流量 $\mathrm{q_m}$，$\mathrm{kg/s}$

% $\quad \mathrm{m^3\cdot s^{-1}}$ and $\quad \mathrm{kg\cdot s^{-1}}$

\begin{equation}
q_m= \rho q_v
\end{equation}

\section{流速}

\bluebold{流速}：单位时间内流体在流动方向上流过的距离，u，$\mathrm{m/s}$

\begin{equation}
u = \frac{q_v}{A}
\end{equation}

\bluebold{圆形管路}

\begin{equation}
u = \frac{4 q_v}{\pi d^2}
\end{equation}

\bluebold{质量流速}：单位时间内流体流过单位截面的质量，G，$\mathrm{kg\cdot m^{-2}\cdot s^{-1}}$

\begin{equation}
G = \mathrm{\rho u}
\end{equation}

\bluebold{流体计算离不开流速和流量的计算}

\chapter{稳态流动与非稳态流动}

\bluebold{稳态流动}：系统中各截面的u、p、$\rho$ 等物理量仅与位置有关，与时间无关。

\bluebold{非稳态流动}：系统中各截面的u、p、$\rho$ 等物理量与时间和位置有关。

\chapter{牛顿黏性定律与流体黏度}

\section{牛顿黏性定律}

\begin{equation}
\tau = \frac{F}{A} = \mu \frac{du}{dy}
\end{equation}

\(\tau = \frac{F}{A} = \mu \frac{du}{dy}\)

\(\tau\)：剪切应力

$\mu$：\(\mu\) 动力黏度

\(\dfrac{du}{dy}\)：速度梯度

详细描述：

\(\boldsymbol{\tau}\):应力张量（矢量/张量形式，后续高阶内容）

\(\tau\)：流体层间剪切应力，内摩擦应力，$\mathrm{N / m^2}$

$\mu$：流体的黏性系数或黏度，$\mu,\ \mathrm{Pa\cdot s}$，在流体里代表\bluebold{动力黏度}。

\(\dfrac{du}{dy}\)：法向速度梯度。

\section{流体黏度}

\begin{equation}
\mu = \frac{T}{\frac{du}{dy}}
\end{equation}

$\mu$：促使流体产生单位速度梯度的剪应力（内摩擦应力），速度梯度up，内摩擦应力up，
黏度总是和速度梯度相关联

黏度是流体的物理性质，$\mu = f(\tau, p)$

$\tau$ up, 液体 $\mu$ down, 气体 $\mu$ up

压力影响可以忽略

\chapter{流动类型和雷诺数}

\section{流动类型}

\bluebold{层流（滞流）}：质点沿轴线方向做规则直线运动，层间质点不混合、不碰撞

\bluebold{湍流（紊流）}：质点轴向运动 + 径向不规则脉动，质点碰撞、混合，大量漩涡，
质点速度大小和方向随时变化

\section{雷诺数}

\begin{equation}
\mathrm{Re} = \frac{d u \rho}{\mu}
\end{equation}

$Re \leq 2000$：滞流 或 湍流

$2000 < Re < 4000$：过渡区，不稳定;

工业上 $Re > 3000$：都按照湍流考虑

$Re \geq 4000$：湍流 或 紊流


\chapter{流体在圆管中的速度分布}

\section{层流流动的速度分布}

同一截面上各点速度不同，抛物线分布

管中心速度最大，管壁处为 0

\begin{equation}
u = \frac{u_\text{max}}{2}
\end{equation}

\section{湍流流动的速度分布}

速度分布比层流均匀，管壁处速度梯度较大

\chapter{流动边界层的概念}

\bluebold{边界层}：管壁附近存在\bluebold{较大速度梯度}的流体层

\bluebold{流体主流区}：边界层外的流体，速度梯度可以忽略，黏性不起作用，可以看成\bluebold{理想流体}

\bluebold{边界层区}：边界层内，速度梯度很大，流动阻力集中在边界层区

湍流边界层也存在\bluebold{层流底层}（层流内层）

\section*{tips}

\bluebold{边界层很少单独出题考试，但是流动阻力、传热、传质理论分析的基础}

\end{document}